% MT048 / explanation — canonical editable content.
考虑一维输入：\(x=-1\)和\(x=1\)属于正类，\(x=0\)属于负类。在原数轴上，正类夹着负类，只用一个阈值把左右分开，无法把两端同时归正、中间归负。


把输入变成\(\phi(x)=x^2\)：原来的两个正点都变成1，负点仍为0。在新轴上，用直线分数\(\phi(x)-0.5\)即可分开。回到原坐标，这条规则是\(x^2-0.5\)，它能给中间与两端不同的类别。这说明特征变换怎样改变可表达的边界。


\StudyImage{CS5489/exam-review/Midterm/assets/answer-kernel-map.png}


\textbf{核技巧使我们能直接计算映射后的内积。} 多维向量的内积是把对应分量相乘再相加，例如\((1,2)\)与\((3,4)\)的内积为\(1\times3+2\times4=11\)。记号\(u^Tv\)中的\(T\)将列向量转成行向量，以完成这次相乘。一维时，内积就是普通乘法，因此本例的核为



\[
k(x,z)=\phi(x)\phi(z)=x^2z^2.
\]


假如\(x=-1,z=2\)，先映射再相乘得到\(1\times4=4\)；直接计算核也得4。两种路径得到同一个内积。对本例这样的小映射，直接构造没有负担；对于极多特征的映射，直接求核值才可能省去大量工作。


SVM可写成按训练样本组合分数的形式：\(f(x)=\sum_i\alpha_i y_i\phi(x_i)^T\phi(x)+b\)。这里\(x_i\)是训练样本，\(y_i\)是其正负标签，\(\alpha_i\)是训练得到的贡献系数，\(b\)是偏置。系数非零、仍参与预测的样本称为支持向量。


把每项内积替换成\(k(x_i,x)\)，就能计算同一个预测分数。采用按样本系数求解的等价训练形式时，训练中的特征运算也只需这些内积，因此不必逐维构造映射后的向量。


\textbf{出处：} 2023A期中 Q9，原卷第6页（10分）


\textbf{答案依据：} 2023A期中 Q9 随卷／配套参考答案；条件与解释经复核。
